%% kapa: exam=1 round=19 year=2008 period=1 form=A subject=civil_law subject_ko=민법 seq=33 number=33
%% source: Q-Net 감정평가사 기출문제 (cst00309) artlSeq=1001962 file=extracted/1008/1001962/2008년 제19회 감정평가사 1차시험 기출문제/1교시.hwp
%% conversion: kordoc 4.24.51 -> tools/gen1.py (auto); review=proofread-sonnet; flags=none
\begin{question}{33}
지역권에 관한 설명으로 틀린 것은? (다툼이 있는 경우에는 판례에 의함)
\begin{choices}
  \item 요역지가 공유관계에 있는 경우 공유자 1인이 지역권을 시효로 취득한 때에는 다른 공유자도 지역권을 취득한다.
  \item 지역권은 요역지와 분리하여 처분할 수 없다.
  \item 乙이 甲소유 토지의 일부를 평온․공연하게 20년 이상 통행하고 있다면, 스스로 통로를 개설하지 않았더라도 통행지역권을 시효취득한다.
  \item 도자기 생산업자가 도자기를 만들기 위해 이웃 토지의 토사를 채취할 목적의 지역권을 취득할 수 없다.
  \item 지역권을 취득한 甲이 乙에게 요역지를 처분하여 소유권이전등기를 마치면 특별한 사정이 없는 한 지역권의 이전등기가 없더라도 지역권이전의 효력이 생긴다.
\end{choices}
\end{question}
\answer{3}
